\chapter{Estado del arte}

Antes de decidir cómo va a funcionar el sistema, revisamos qué herramientas y trabajos
existen para el análisis automatizado del FST. La búsqueda se centró en tres áreas:
software para análisis conductual, métodos de seguimiento de animales con visión por
computadora, y enfoques de clasificación de conducta con aprendizaje profundo.

\section{Herramientas comerciales de análisis conductual por video}

\subsection{EthoVision XT}
EthoVision XT (Noldus Information Technology) es el software más usado en laboratorios
de neurociencia conductual para seguimiento automático de animales en video
\cite{noldus_ethovision}. Detecta al animal en cada cuadro por sustracción de fondo y
umbralización, y registra posición, velocidad y tiempo de inmovilidad.

Para el FST, EthoVision puede detectar si el animal se mueve o no, pero no distingue
entre los tres tipos de conducta (nado, inmovilidad, escalamiento) que importan para el
análisis farmacológico \cite{della_valle2025}. Además, necesita condiciones de
iluminación bastante controladas para funcionar bien, y su licencia comercial tiene un
costo que no todos los laboratorios pueden cubrir.

\subsection{ANY-maze}
ANY-maze (Stoelting Co.) tiene un uso similar al de EthoVision y una interfaz más
accesible, pero comparte las mismas limitaciones de fondo: requiere condiciones visuales
controladas y no clasifica las subconductas del FST por separado.

\section{Herramientas de código abierto para análisis del FST}

\subsection{ezTrack}
ezTrack \cite{pennington2019_eztrack} es una herramienta de código abierto en Python,
pensada para investigadores sin experiencia en programación. Corre sobre Jupyter
Notebooks y cuantifica el movimiento comparando cuadros consecutivos.

Sirve para experimentos simples, pero solo hace clasificación binaria (se mueve/no se
mueve) y asume condiciones de grabación razonablemente controladas. No fue diseñada
para vista superior con varios cilindros al mismo tiempo.

\subsection{DBscorer}
DBscorer \cite{nandi2021_dbscorer} es la herramienta de código abierto más específica
para el FST y el test de suspensión de cola (TST). La desarrolló el Indian Institute of
Science en MATLAB y tiene una interfaz gráfica bastante intuitiva. Cuantifica la
inmovilidad comparando el tamaño de la región del animal entre cuadros consecutivos.

Es un paso adelante frente al análisis manual, pero tiene limitaciones concretas para
este proyecto: es sensible a cambios de iluminación y reflejos del agua, solo distingue
móvil/inmóvil sin separar nado de escalamiento, está diseñada para un solo animal a la
vez en perspectiva lateral, y no tiene interfaz web.

\section{Enfoques de aprendizaje profundo para el FST}

\subsection{Red dual sin umbral (DSAAN)}
Meng et al.\ \cite{meng2024_dsaan} propusieron DSAAN, un método con dos ramas paralelas
que procesan apariencia y movimiento del animal por separado y luego los combinan para
clasificar el estado cuadro a cuadro. Con apenas 94 imágenes etiquetadas alcanzó un
88.7\,\% de precisión, lo cual es relevante dado el tamaño de dataset que tendríamos
disponible.

El método se validó en condiciones controladas; los autores no reportan qué pasa con
variaciones de iluminación o reflejos del agua.

\subsection{Red convolucional 3D-RCNN (Della Valle et al., 2025)}
El trabajo más reciente que encontramos sobre clasificación automática en el FST.
Della Valle et al.\ \cite{della_valle2025} entrenaron un modelo basado en redes
residuales 3D que procesa segmentos de video para capturar el movimiento en el tiempo.
Es uno de los pocos sistemas que distingue las tres conductas; los propios autores
señalan que la mayoría de sistemas anteriores solo diferencian entre inmovilidad y
conductas activas en general.

También fue validado en condiciones controladas.

\section{Estimación de pose y seguimiento multianimales}

\subsection{DeepLabCut}
DeepLabCut \cite{mathis2018_dlc} es la herramienta estándar para seguimiento de pose
en animales sin marcadores físicos. Usa redes convolucionales para detectar puntos clave
del cuerpo (nariz, patas, cola) definidos por el usuario. Su versión multianimales,
maDLC \cite{lauer2022_madlc}, permite trabajar con varios individuos al mismo tiempo.

Sturman et al.\ \cite{sturman2020_dlc} mostraron que un clasificador entrenado sobre
la representación esquelética de DeepLabCut puede eliminar la variabilidad entre
evaluadores en el FST. El precio es el esfuerzo de etiquetado: hay que anotar
manualmente puntos clave en decenas de cuadros antes de poder entrenar el modelo.

\subsection{YOLO-Behaviour}
Chan et al.\ \cite{chan2025_yolo} presentaron un sistema que combina YOLOv8 para
detectar al animal con un clasificador de conducta que actúa sobre la región recortada.
Separar detección de clasificación tiene una ventaja práctica: el clasificador trabaja
sobre una imagen más limpia, sin distracciones del fondo. Ese enfoque de dos etapas es
el que se retoma en este proyecto.

\subsection{YOLOv8 para detección en investigación conductual}
Stadler et al.\ \cite{stadler2024_yolov8} evaluaron YOLOv8 para detección de animales
en estudios de comportamiento. Con entre 100 y 350 imágenes etiquetadas lograron
detección casi perfecta en fondos consistentes. También encontraron que el modelo no
generaliza bien cuando el fondo cambia mucho entre condiciones, lo cual es un punto a
tener en cuenta al entrenar el sistema.

\section{Limitaciones comunes y contexto del proyecto}
Todos los trabajos revisados se desarrollaron y validaron bajo condiciones bastante
controladas: iluminación uniforme, fondo contrastante, sin reflejos problemáticos. Las
herramientas basadas en umbralización (DBscorer, ezTrack) son especialmente sensibles
a eso. Los métodos de aprendizaje profundo son más robustos en principio, pero su
rendimiento depende de que el conjunto de entrenamiento represente las condiciones
reales de uso.

Los videos del laboratorio colaborador tienen características que complican las cosas:
grabación con cámara de celular, variaciones de iluminación entre sesiones, y ratas
de pelaje claro con poco contraste frente al cilindro. El sistema propuesto no garantiza
robustez ante todas esas condiciones desde el inicio; parte del trabajo de desarrollo
será evaluar qué tan bien funciona con los videos disponibles y definir los requisitos
mínimos de calidad de entrada. La hipótesis es que los videos son suficientemente buenos
para que el sistema funcione, pero eso se confirma o se ajusta durante las pruebas.

\section{Tabla comparativa}
La Tabla~\ref{tab:comparativa_herramientas} resume los trabajos revisados frente a los
criterios más relevantes para el proyecto.

\begin{table}[ht]
\centering
\caption{Comparativa de herramientas y trabajos relacionados con el análisis del FST.}
\label{tab:comparativa_herramientas}
\small
\begin{tabular}{lccccc}
\hline
\textbf{Sistema} & \textbf{Espec. FST} & \textbf{3 conductas} & \textbf{DL} & \textbf{Web} & \textbf{Abierto} \\
\hline
EthoVision XT      & Parc. & No    & No  & No  & No  \\
ANY-maze           & Parc. & No    & No  & No  & No  \\
ezTrack            & Parc. & No    & No  & No  & Sí  \\
DBscorer           & Sí    & No    & No  & No  & Sí  \\
DSAAN              & Sí    & No    & Sí  & No  & No  \\
3D-RCNN            & Sí    & Sí    & Sí  & No  & No  \\
DeepLabCut         & Parc. & Parc. & Sí  & No  & Sí  \\
YOLO-Behaviour     & No    & ---   & Sí  & No  & Sí  \\
\hline
\textbf{Propuesta} & \textbf{Sí} & \textbf{Sí} & \textbf{Sí} & \textbf{Sí} & \textbf{Sí} \\
\hline
\end{tabular}
\end{table}

Ningún sistema existente cubre los cinco criterios al mismo tiempo. Los comerciales no
son abiertos y no clasifican las tres conductas. Los de código abierto no usan
aprendizaje profundo ni tienen interfaz web. Los de aprendizaje profundo más recientes
sí clasifican las tres conductas pero no tienen plataforma web. El sistema propuesto
busca cubrir los cinco, con la salvedad de que el desempeño real dependerá de la
calidad de los videos con los que se trabaje.